% !PRINTED none
\weberpage{none}
% [source note: this Abschnitt opener has no printed folio; sequence position corresponds to inferred page 705.]

\begin{center}
\gesperrt{Vierundzwanzigster Abschnitt.}

\gesperrt{Classenzahl der Kreistheilungskörper.}

\rule{4em}{0.4pt}
\end{center}

\begin{center}
\S.\,193.

\gesperrt{Classenzahldarstellung im Kreistheilungskörper $\Omega_m$.}
\end{center}

Die allgemeinen Resultate über die Classenzahl $h$ sollen
nun angewandt werden auf den vollen Kreistheilungskörper $\Omega_m$
unter der bisherigen Voraussetzung, dass $m$ eine Potenz einer
Primzahl $q$ sei, und dass also
\weq{1}{m = q^{\varkappa}, \quad \varphi(m) = q^{\varkappa-1}(q-1) = \mu}
gesetzt sei.

Im einundzwanzigsten Abschnitte (\S.\,169, 171) haben wir
gesehen, dass in $q$ nur ein Primideal $1^{\text{sten}}$ Grades aufgeht, und
dass eine zum Exponenten $f$ \gesperrt{gehörige}, von $q$ verschiedene
Primzahl $p$ in $e$ verschiedene Primideale $f^{\text{ten}}$ Grades zerfällt.
Darin bedeutet $f$ den kleinsten positiven Exponenten, für den
\weq{2}{p^f \equiv 1\ (\mathrm{mod}\ m)}
ist, und $e$ ist durch die Gleichung
\weq{3}{\mu = ef}
bestimmt.

Die am Ende des vorigen Paragraphen durch die Formel (6)
bestimmte Function $\varPhi(s)$ erhält daher hier den Ausdruck:
\weq{4}{\varPhi(s) = \frac{1}{1-q^{-s}}\prod\frac{1}{(1-p^{-sf})^e},}
und darin erstreckt sich das Productzeichen $\varPi$ auf alle Primzahlen
 $p$, die von $q$ verschieden sind. Darin ist $s$ eine Variable,
die immer grösser als 1 ist, die sich schliesslich der Grenze 1
nähern soll.
